\def\CTeXPreproc{Created by ctex v0.2.12, don't edit!}
\documentclass[12pt]{article}
\usepackage{}

\usepackage{epsfig}
\usepackage{latexsym}
%\usepackage{algorithmic}
\usepackage{caption}
\usepackage{amsfonts}
\usepackage{amssymb}
\usepackage{mathrsfs}
\usepackage{amsmath}
\usepackage{enumerate}
\usepackage{graphics}
\usepackage{float}
\usepackage{pict2e}


\usepackage{color}
\def \red {\textcolor{red} }
\def \bl  {\textcolor{blue} }
\def \gr  {\textcolor{green} }

\setlength{\textheight}{8.5in} \setlength{\textwidth}{6.2in}
\setlength{\oddsidemargin}{0in} \setlength{\parindent}{1em}

\makeatletter
\newcommand{\rmnum}[1]{\romannumeral #1}
\newcommand{\Rmnum}[1]{\expandafter\@slowromancap\romannumeral #1@}

\newcommand{\rr}{Res_{k+1}}
\makeatother

\begin{document}

%\pagestyle{myheadings} \markright{\sc
%$(k\mathcal{C},\mathcal{C}_{d})$-decomposition of Graphs \hfill}
%\thispagestyle{empty}

\newtheorem{theorem}{Theorem}
\newtheorem{observation}[theorem]{Observation}
\newtheorem{corollary}[theorem]{Corollary}
\newtheorem{algorithm}[theorem]{Algorithm}
\newtheorem{definition}[theorem]{Definition}
\newtheorem{guess}[theorem]{Conjecture}
\newtheorem{claim}[theorem]{Claim}
\newtheorem{problem}[theorem]{Problem}
\newtheorem{question}[theorem]{Question}
\newtheorem{lemma}[theorem]{Lemma}
\newtheorem{proposition}[theorem]{Proposition}
\newtheorem{fact}[theorem]{Fact}

%\renewcommand{\algorithmicrequire}{\textbf{Input}}

%\renewcommand{\algorithmicensure}{\textbf{Output}}
\makeatletter
  \newcommand\figcaption{\def\@captype{figure}\caption}
  \newcommand\tabcaption{\def\@captype{table}\caption}
\makeatother


\newtheorem{acknowledgement}[theorem]{Acknowledgement}

\newtheorem{axiom}[theorem]{Axiom}
\newtheorem{case}[theorem]{Case}
\newtheorem{conclusion}[theorem]{Conclusion}
\newtheorem{condition}[theorem]{Condition}
\newtheorem{conjecture}[theorem]{Conjecture}
\newtheorem{criterion}[theorem]{Criterion}
\newtheorem{example}[theorem]{Example}
\newtheorem{exercise}[theorem]{Exercise}
\newtheorem{notation}{Notation}
%% \newtheorem{remark}[theorem]{Remark}
\newtheorem{solution}[theorem]{Solution}
\newtheorem{summary}[theorem]{Summary}
%% \newenvironment{proof}[1][Proof]{\noindent\textbf{#1.} }{\ \rule{0.5em}{0.5em}}

\newenvironment{proof}{\noindent {\bf
Proof.}}{\rule{3mm}{3mm}\par\medskip}
\newcommand{\remark}{\medskip\par\noindent {\bf Remark.~~}}
\newcommand{\pp}{{\it p.}}
\newcommand{\de}{\em}
%% \newcommand{\mod}{ \ {\rm mod} \ }
\newcommand{\mad}{\rm mad}
\newcommand{\qf}{Q({\cal F},s)}
\newcommand{\qff}{Q({\cal F}',s)}
\newcommand{\qfff}{Q({\cal F}'',s)}
\newcommand{\f}{{\cal F}}
\newcommand{\ff}{{\cal F}'}
\newcommand{\fff}{{\cal F}''}
\newcommand{\fs}{{\cal F},s}
\newcommand{\s}{\mathcal{S}}
\newcommand{\G}{\Gamma}
\newcommand{\g}{\gamma}
\newcommand{\wrt}{with respect to }
\newcommand {\nk}{ Nim$_{\rm{k}} $  }

\newcommand{\q}{\uppercase\expandafter{\romannumeral1}}
\newcommand{\qq}{\uppercase\expandafter{\romannumeral2}}
\newcommand{\qqq}{\uppercase\expandafter{\romannumeral3}}
\newcommand{\qqqq}{\uppercase\expandafter{\romannumeral4}}
\newcommand{\qqqqq}{\uppercase\expandafter{\romannumeral5}}
\newcommand{\qqqqqq}{\uppercase\expandafter{\romannumeral6}}

\newcommand{\qed}{\hfill\rule{0.5em}{0.809em}}

\newcommand{\var}{\vartriangle}


\title{{\large \bf Multiple list colouring of  planar graphs}}
%\title{{The properties of the hypercube-like}}


\author{   Xuding Zhu\thanks{Department of Mathematics, Zhejiang Normal University,  China.  E-mail: xudingzhu@gmail.com. Grant Number: NSFC 11571319.},}

\date{}
\maketitle

\begin{abstract}
This paper proves that for each positive integer $m$, there is a planar graph $G$ which is not $(4m+\lfloor \frac{2m-1}{9}\rfloor,m)$-choosable. Then we pose some conjectures concerning multiple list colouring of planar graphs.  
\end{abstract}

\vspace{3mm}\textbf{Keywords:} fractional choice number,  multiple list colouring,  planar graph.


\section{Introduction}

A $b$-fold colouring of a graph $G$ is a mapping $\phi$ which assigns to each vertex $v$ of $G$ a set $\phi(v)$ of $b$ colours, so that adjacent vertices receive disjoint colour sets. An $(a,b)$-colouring of $G$ is a $b$-fold colouring $\phi$ of $G$ such that $\phi(v) \subseteq \{1,2,\ldots, a\}$ for each vertex $v$. The {\em fractional chromatic number} of $G$ is
 $$\chi_f(G)= \inf\{\frac ab: G \ \text{ is $(a,b)$-colourable } \}.$$
An $a$-list assignment of $G$ is a mapping $L$ which assigns to each vertex $v$ a set $L(v)$ of $a$ permissible colours. A $b$-fold $L$-colouring of $G$ is a $b$-fold colouring $\phi$ of $G$ such that $\phi(v) \subseteq L(v)$ for each vertex $v$. We say $G$ is $(a,b)$-choosable if for any $a$-list assignment $L$ of $G$, there is a $b$-fold $L$-colouring of $G$. The {\em fractional choice number} of $G$ is $$ch_f(G)= \inf\{\frac ab:  G \ \text{is $(a,b)$-choosable }\}.$$ It was proved by Alon, Tuza and Voigt \cite{ATV1995} that for any finite graph $G$, $\chi_f(G)=ch_f(G)$ and moreover the infimum in the definition of $ch_f(G)$ is attained and hence can be replaced by minimum. This implies that if $G$ is $(a,b)$-colourable, then for some integer $m$, $G$ is $(am,bm)$-choosable. As every planar graph is $4$-colourable (which is equivalent to $(4,1)$-colourable), we know for each planar graph $G$, there is an integer $m$ such that $G$ is $(4m,m)$-choosable. However, the integer $m$ depends on $G$. 
We prove in this paper that there is no integer $m$ so that every planar graph is $(4m,m)$-choosable.   

\begin{theorem}
	\label{main}
	For each positive integer $m$, there is a planar graph $G$ which is not $(4m+\lfloor \frac{2m-1}{9}\rfloor,m)$-choosable.
\end{theorem}

The $m=1$ case of Theorem \ref{main} is equivalent to say that there are non-$4$-choosable planar graphs, which was  proved by   Voigt \cite{Voigt1993}. A smaller non-$4$-choosable planar graph was constructed \cite{Mir1996}, a $3$-colourable non-$4$-choosable planar graph was constructed by Gutner   \cite{Gut1996,VW1997} and it was prove in \cite{Gut1996} that it is NP-complete to decide if a given planar graph is $4$-choosable.


% The current known smallest non-$4$-choosable planar graph has $63$ vertices. 

 

\section{The proof of Theorem \ref{main}}

In this section, $m$ is a fixed positive integer. Let $k= \lfloor \frac{2m-1}{9}\rfloor$. We shall construct a planar graph 
$G$ which is not $(4m+k,m)$-choosable.


\begin{figure}[hh!]
	\label{fig2}
	\centering
	\includegraphics[width=140mm]{G.pdf}
	\label{fig1}
	\caption{The graph $G$}
\end{figure}

\begin{lemma}
	\label{lem1}
	Let $G$ be the graph as shown in Figure 1. There is a list assignment $L$
	of $G$ for which the following hold:
	\begin{enumerate}
		\item $|L(s)|=4m+k$ for each vertex $s$, except that  $|L(u)|=|L(u')|=m$.
		\item There is no $m$-fold $L$-colouring of $G$.
	\end{enumerate} 
\end{lemma}
\begin{proof}
	Let $A, B, C,D$ be pairwise disjoint sets of colours such that $|A|=|B|=m$ and $|C|=|D|=2m+k$. Let $X, X'$ be disjoint subsets of $C$ with $|X|=|X'|=m$. Define list assignment $L$ as follows:
	\begin{itemize}
		\item $L(u)=A$ and $L(u')=B$.
		\item $L(v)=L(w)=L(t)=L(t')=A \cup B \cup C$.
		\item $L(x)=L(a)=X \cup A \cup D$, and $L(x')=L(a')=X' \cup A \cup D$.
		\item $L(y)=L(b)=X \cup B \cup D$, and $L(y')=L(b')=X' \cup B \cup D$.
		\item $L(z)=L(c)= L(z')=L(c')=A \cup B \cup D$.
	\end{itemize}
	We shall show that there is no $m$-fold $L$-colouring of $G$.
	Assume to the contrary that $\phi$ is an $m$-fold $L$-colouring of $G$.
	Then $\phi(u)=A$ and $\phi(u')=B$ and $\phi(v), \phi(w)$ are two disjoint $m$-subsets of $C$.
	
	Let $S=G-(\phi(v) \cup \phi(w))$. We have $|S|=k$.
	As $X, X'$ are  disjoint subsets of $C$, we have
	$$|X \cap S|+|X' \cap S| \le |S|=k.$$
	By symmetry between the two triangles $(u,v,w)$ and $(u',v,w)$, we may assume that 
	$$|X \cap S| \le |X' \cap S|, \ \text{and hence} \ |X \cap S| \le k/2.$$
	(In case $|X' \cap S| \le \frac k2$, then we consider the subgraph contained in the triangle $(u',v,w)$).
	Thus
	$$|X \cap \phi(v)| +|X \cap \phi(w)|=|X \cap (\phi(v) \cup \phi(w))| \ge m- \frac k2.$$
	By symmetry between triangles $(u,v,t)$ and $(u,w,t)$, we may assume that 
	$$|X \cap \phi(v)| \ge |X \cap \phi(w)|, \ \text{and hence } \ |X \cap \phi(v)| \ge \frac m2 - \frac k4.$$
		(In case $|X \cap \phi(w)| \le \frac m2 - \frac k4$, then we consider the subgraph contained in the triangle $(u,w,t)$).
		
	Let $T=X-\phi(v)$. We have $$|T|=|X|-|X \cap \phi(v)| \le \frac m2 + \frac k4.$$
	Let $R=B - \phi(t)$. As $\phi(t)$ is disjoint from $\phi(u) \cup \phi(v) \phi(w)$,
	we know that $\phi(t) \subseteq B \cup S$. Hence
	$$|R|\le |S|= k.$$
By deleting the colours used by the neighbours of $a,b,c$, respectively, we have
	\begin{itemize}
		\item $\phi(a) \subseteq D \cup T$,
		\item $\phi(b) \subseteq D \cup R$,
		\item $\phi(c) \subseteq D \cup T \cup R$.
	\end{itemize}
	As $\phi(a),\phi(b),\phi(c)$ are pairwise disjoint, we have
	$$3m = |\phi(a) \cup \phi(b) \cup \phi(c)| \le |D|+|T|+|R| 
	\le (2m+k)+\left(\frac m2 + \frac k4\right) + k = \frac{5m}{2}+ \frac{9k}{4} < 3m,$$
	a contradiction.
\end{proof}
 
Let $p={4m+k \choose m, m,2m+k}$, and let $H$ be obtained from the disjoint union of $p$ copies of $G$ by identifying all the copies of $u$ into a single vertex (also named as $u$) and all the copies of $u'$ into a single vertex (also named as $u'$), and then add an edge connecting $u$ and $u'$. It is obvious that $H$ is a planar graph. 

Now we show that   $H$ is not $(4m+k,m)$-choosable.  Let $Z$ be a set of $4m+k$ colours.
Let $L(u)=L(u')=Z$. There are $p$ possible $m$-fold $L$-colourings of $u$ and $u'$. Each such a colouring $\phi$ corresponds to one copy of $G$. In that copy of $G$, define the list assignment as in the proof of Lemma \ref{lem1}, by replacing $A$ with $\phi(u)$ and $B$ with $\phi(u')$.  Now Lemma \ref{lem1} implies that no $m$-fold colouring of $u$ and $u'$ can be extended to an $m$-fold $L$-colouring of $H$. This completes the proof of Theorem \ref{main}.


\section{Some open problems}

Thomassen proved that every planar graph is $5$-choosable \cite{Tho1994}. The proof can be easily adopted to show that for any positive integer $m$, every planar graph is $(5m,m)$-choosable. Given a positive integer $m$, let $a(m)$ be the minimum integer such that every planar graph is $(a(m),m)$-choosable.
Combining Thomassen's result and Theorem \ref{main}, we have $$4m+\lfloor \frac{2m-1}{9}\rfloor+1 \le a(m) \le 5m.$$
 For $m=1$, the upper bound and the lower bound coincide. So $a(1)=5$. As $m$ becomes bigger, the gap between the upper and lower bounds increases.
 A natural question is what is the exact value of $a(m)$. We conjecture that the upper bound is not always tight.
 
   

\begin{conjecture}
	\label{conj1}
	There is a constant integer $m$ such that every planar graph is $(5m-1,m)$-choosable.
\end{conjecture}

It is proved recently in \cite{HZ2016} that every planar graph is   $1$-defective
$(9,2)$-paintable, which implies that every planar graph is $1$-defective $(9,2)$-choosable (i.e., if each vertex has $9$ permissible colours, then there is a
$2$-fold colouring of the vertices of $G$ with permissible colours so that each colour class induces a graph of maximum degree at most $1$.) 
The following conjecture, which is stronger than Conjecture \ref{conj1}, asserts that the $1$-defective can be replaced by $0$-defective.

\begin{conjecture}
	\label{conj2}
	Every planar graph is $(9,2)$-choosable.
\end{conjecture}


 

It follows from the Four Colour Theorem  that for any integer $m$, every planar graph is $(4m,m)$-colourable. Without using the Four Colour Theorem, it is proved very recently by Cranston and Rabern  \cite{CR2016}  that every planar graph is $(9,2)$-colourable (an earlier result in \cite{HRS1973} shows that every planar graph $G$ is $(5m-1,m)$-colourable with $m = |V(G)|+1$).


\bibliographystyle{unsrt}
\begin{thebibliography}{99}
	
\bibitem{ATV1995}	N. Alon, Zs. Tuza, and M. Voigt, {\em Choosability and fractional chromatic numbers}, Discrete Math.,	165/166:31--38, 1997. Graphs and combinatorics (Marseille, 1995).
	
	\bibitem{CR2016} D. Cranston and L. Rabern, {Planar graphs are $9/2$-colorable}, manuscript, 2015.
	
\bibitem{Gut1996} S.	Gutner,  
{\em	The complexity of planar graph choosability},
	Discrete Math. 159 (1996), no. 1-3, 119--130. 
	
	\bibitem{HZ2016} M. Han and X. Zhu, {\em Every planar graph is $1$-defective $(9,2)$-paintable}, manuscript, 2016.
	
\bibitem{HRS1973}	A.J.W. Hilton, R.~Rado, and  S.H. Scott, {\em   A $(<5)$-colour theorem for planar graphs},
  Bull. London Math. Soc., 5:302--306, 1973.
	
	
	\bibitem{Mir1996} M. Mirzakhani,{\em  
		A small non-4-choosable planar graph},
	Bull. Inst. Combin. Appl. 17 (1996), 15--18. 

 
\bibitem{Tho1994} C. Thomassen, {\em Every planar graph is $5$-choosable}, J. Combin. Theory B, 64(1994),101--107.   

\bibitem{Voigt1993} M. Voigt {\em List colourings of planar graphs}, Disc. Math., 120(1993), 215--219. 

\bibitem{VW1997} M.  Voigt and B. Wirth, {\em  On 3-colorable non-4-choosable planar graphs}, J. Graph Theory 24 (1997), no. 3, 233--235.

\end{thebibliography}

\end{document} 

%

Let $G_0$ be the plane graph as shown in Fig. \ref{fig1}. 

\begin{figure}[hh!]
	\label{fig1}
	\centering
	\includegraphics[width=80mm]{G0.pdf}
	\caption{The graph $G_0$}
\end{figure}


Let $G_1$ be the plane graph obtained from the disjoint union of ${8 \choose 2}^2$ copies of $G_0$, by adding two vertices $u$ and $w$, where $u$ is adjacent to all copies of $v_1,v_2, v_3,v_4$ and $u_1,u_2,u_3$, and $w$ is adjacent to all copies of $v_1,v_2,v_3,v_4$ and $w_1,w_2,w_3$. (Vertices in the $j$th copy of $G_0$ are named as $v_i^j, u_i^j$ and $w_i^j$.)



\begin{figure}[hh!]
	\label{fig2}
	\centering
	\includegraphics[width=140mm]{G1.pdf}
	\caption{The graph $G_1$}
\end{figure}


Let $L$ be the $8$-list assignment of $G_1$ defined as follows: $L(u)=\{a,b,c,d,e,f,g,h\}$ and $L(w)=\{a',b',c',d',e',f',g',h'\}$. 
There are $k={8 \choose 2}^2$ different $2$-fold $L$-colourings of the two vertices $u$
and $w$. Label these $k$ colourings of $u$ and $w$ as $f_1,f_2, \ldots, f_k$.
For $1 \le j \le k$, the list assignment of the $j$th copy of $G_0$ is given as follows:

\begin{itemize}
	\item $L(v_1^j)=L(v_2^j)=L(v_3^j)=L(v_4^j)=\{1,2,3,4\} \cup f_j(u) \cup f_j(w)$. 
	\item $L(u_1^j)=L(u_2^j)=L(u_3^j)=\{1,2,3,4,5,6\} \cup f_j(u)$.
	\item $L(w_1^j)=L(w_2^j)=L(w_3^j)=\{1,2,3,4,5,6\} \cup f_j(w)$.
\end{itemize}

\begin{lemma}
	For any $2$-fold $L$-colouring $\phi$ of $G_1$, there is an index $j$ such that \begin{itemize}
		\item $\phi((v_1^j), L(v_2^j),L(v_3^j),L(v_4^j) \subseteq \{1,2,3,4\})$.
		\item $\phi(u_1^j)=\phi(u_2^j)=\phi(u_3^j)=\phi(w_1^j)=\phi(w_2^j)=\phi(w_3^j)=\{ 5,6\}$.
	\end{itemize}  
	
\end{lemma}
\begin{proof}
	Assume $\phi$ is a $2$-fold $L$-colouring of $G_1$. The restriction of $\phi$ to $u,w$ equals $f_j$ for some $j \in \{1,2,\ldots,k\}$.
	This implies that  $\phi((v_1^j), L(v_2^j),L(v_3^j),L(v_4^j) \subseteq \{1,2,3,4\})$, and $\phi(u_1^j), \phi(u_2^j), \phi(u_3^j), \phi(w_1^j), \phi(w_2^j), \phi(w_3^j)\subseteq \{1,2,3,4, 5,6\}$.  This implies that $\phi((v_1^j) \cup L(v_2^j) = \{1,2,3,4\}$ and hence $\phi(u_1^j)=\phi(w_1^j)=\{5,6\}$.
	Similarly, we have $\phi(u_2^j)=\phi(w_2^j)=\phi(u_3^j)=\phi(w_3^j)=\{5,6\}$.
\end{proof}

Consider the graph $G_2$ shown in Figure \ref{fig3}. Let $L$ be the list assignment of $G_2$ defined as follows:

\begin{itemize}
	\item  $L(u_i)=L(w_i)=\{5,6\}$ for $i=1,2,3$.
	\item $L(v_i)=\{1,2,3,4\}$, $L(x_i)=\{1,2,3,4,a,b,c,d\}$ for $i=1,2,3,4$.
	\item $L(y_1)=L(y'_1)=\{3,4,5,6,a,b,c,d\}$.
	\item $L(z_1)=L(z'_1)=\{1,2,5,6,a,b,c,d\}$.
	\item $L(y_2)=L(y'_2)=\{2,4,5,6,a,b,c,d\}$.
	\item $L(z_2)=L(z'_2)=\{1,3,5,6,a,b,c,d\}$.
	\item $L(y_3)=L(y'_3)=\{2,3,5,6,a,b,c,d\}$.
	\item $L(z_1)=L(z'_1)=\{1,4,5,6,a,b,c,d\}$.
\end{itemize}

\begin{figure}[hh!]
	\label{fig2}
	\centering
	\includegraphics[width=140mm]{G2.pdf}
	\caption{The graph $G_2$}
\end{figure}

\begin{lemma}
	There is no $2$-fold $L$-colouring of $G_2$.
\end{lemma}
\begin{proof}
	Assume $\phi$ is a $2$-fold $L$-colouring of $G_2$. Then $\phi(u_i)=\phi(w_i)=\{5,6\}$ for $i=1,2,3$, and $\phi(v_1) =\phi(v_3) \subset \{1,2,3,4\}$, $\phi(v_2)=\phi(v_4) = \{1,2,3,4\} - \phi(v_1)$.
	By symmetry, we may assume that $1 \in \phi(v_2)$. There are three cases.
	
	\noindent
	{\bf Case 1} $\phi(v_2)=\{1,2\}$ and $\phi(v_1)=\{3,4\}$. In this case, the colours
	permissible for $x_1,y_1,z_1$ are $\{a,b,c,d\}$, and hence cannot be coloured.
	
	\noindent
	{\bf Case 2} $\phi(v_2)=\{1,3\}$ and $\phi(v_1)=\{2,4\}$. In this case, the colours
	permissible for $x_2,y_2,z_2$ are $\{a,b,c,d\}$, and hence cannot be coloured.
	
	\noindent
	{\bf Case 3} $\phi(v_2)=\{1,3\}$ and $\phi(v_1)=\{2,3\}$. In this case, the colours
	permissible for $x_3,y_3,z_3$ are $\{a,b,c,d\}$, and hence cannot be coloured.
\end{proof}






Let $p = {4m \choose m}^2$ and $k = {2m \choose m}$. 
The graph $G$ is as depicted in Figure \ref{fig1}  (in each triangle induced by $\{x_i, z_{i,j}, z_{i,j+1}\}$, a triangle as indicated at left bottom is insered).  


\begin{figure}
	\setlength{\unitlength}{1mm}
	$$\begin{picture}(100,110)
	\thicklines
	\put(50,110){\circle*{1.5}}
	\put(49,111){$u$}
	\linethickness{.3mm}
	\put(50,110){\line(-2.5,-1){60}}
	\put(50,110){\line(-1.6,-1){38}}
	\put(-10,86){\circle*{1.5}}
	\put(-12,88){$x_1$}
	\put(-10,86){\line(1,0){22}}
	\put(12,86){\circle*{1.5}}
	\put(9,88){$y_1$}
	\put(30,90){$\cdots$}
	
	\put(50,110){\line(-0.4,-1){9.5}}
	\put(50,110){\line(0.4,-1){9.5}}
	\put(40.5,86){\circle*{1.5}}
	\put(37,88){$x_i$}
	\put(40.5,86){\line(1,0){19.5}}
	\put(59.5,86){\circle*{1.5}}
	\put(59.5,88){$y_i$}
	\put(69,90){$\cdots$}
	
	\put(50,110){\circle*{1.5}}
	\put(50,110){\line(2.5,-1){60}}
	\put(50,110){\line(1.6,-1){38}}
	\put(88,86){\circle*{1.5}}
	\put(89,88){$x_p$}
	\put(88,86){\line(1,0){22}}
	\put(110,86){\circle*{1.5}}
	\put(110,88){$y_p$}
	
	\put(50,30){\circle*{1.5}}
	\put(49,27){$z$}
	\put(-10,86){\line(1,-0.93){60}}
	\put(12,86){\line(0.8,-1.17){38}}
	\put(40.5,86){\line(0.1,-0.59){9.5}}
	\put(59.5,86){\line(-0.1,-0.59){9.5}}
	\put(110,86){\line(-1,-0.93){60}}
	\put(88,86){\line(-0.8,-1.17){38}}
	
	\put(8,80){\circle*{1.5}}
	\put(4,83){$z_{1,1}$}
	\put(14,73){\circle*{1.5}}
	\put(36,46.5){\circle*{1.5}}
	\put(37,47){$z_{1,k}$}
	
	\put(8,80){\line(1,-1.19){42}}
	\put(-10,86){\line(1,-0.32){18}}
	\put(12,86){\line(-1,-1.4){4.5}}
	\put(-10,86){\line(1,-0.55){24}}
	\put(12,86){\line(1,-6){2.2}}
	\put(6,72){$\ddots$}
	\put(16,72){$\ddots$}
	\put(-10,86){\line(1,-0.865){46}}
	\put(12,86){\line(1,-1.64){24}}
	
	\put(92,80){\circle*{1.5}}
	\put(93,83){$z_{p,1}$}
	\put(86,73){\circle*{1.5}}
	\put(64,46.5){\circle*{1.5}}
	\put(64,44.5){$z_{p,k}$}
	
	\put(92,80){\line(-1,-1.19){42}}
	\put(110,86){\line(-1,-0.32){18}}
	\put(88,86){\line(1,-1.4){4.5}}
	\put(110,86){\line(-1,-0.55){24}}
	\put(88,86){\line(-1,-6){2.2}}
	\put(84,75){\circle*{0.5}}
	\put(83,74){\circle*{0.5}}
	\put(82,73){\circle*{0.5}}
	\put(94,75){\circle*{0.5}}
	\put(93,74){\circle*{0.5}}
	\put(92,73){\circle*{0.5}}
	\put(110,86){\line(-1,-0.865){46}}
	\put(88,86){\line(-1,-1.64){24}}
	
	
	\put(50,80){\circle*{1.5}}
	\put(50,80){\line(0,-1){50}}
	\put(48,82){$z_{i,1}$}
	\put(47,77){$\vdots$}
	\put(52,77){$\vdots$}
	\put(50,70){\circle*{1.5}}
	\put(52,70){$z_{i,j}$}
	\put(50,63){\circle*{1.5}}
	\put(52,63){$z_{i,j+1}$}
	\put(47,57){$\vdots$}
	\put(52,57){$\vdots$}
	
	\put(40.5,86){\line(1,-0.6){9.5}}
	\put(59.5,86){\line(-1,-0.6){9.5}}
	\put(40.5,86){\line(1,-1.7){9.5}}
	\put(59.5,86){\line(-1,-1.7){9.5}}
	\put(40.5,86){\line(1,-2.5){9.5}}
	\put(59.5,86){\line(-1,-2.5){9.5}}
	
	\linethickness{.7mm}
	\put(22,43){\vector(1,1.1){25.4}}
	
	\put(-10,50){\circle*{1.5}}
	\put(-12,52){$x_i$}
	\linethickness{.3mm}
	\put(-10,50){\line(1,-0.29){35}}
	\put(-10,50){\line(1,-1.43){35}}
	\put(-10,50){\line(1,-0.54){18}}
	\put(-10,50){\line(1,-0.95){23}}
	
	\put(25,40){\circle*{1.5}}
	\put(27,38){$z_{i, j}$}
	\put(25,40){\line(0,-1){40}}
	\put(25,0){\circle*{1.5}}
	\put(27,1){$z_{i, j+1}$}
	
	\put(8,40){\circle*{1.5}}
	\put(6,42){$a_{i, j}$}
	\put(8,40){\line(1,-0.01){17}}
	\put(8,40){\line(1,-2.3){5.5}}
	\put(8,40){\line(1,-0.01){17}}
	\put(8,40){\line(1,-0.8){9}}
	
	\put(13,28){\circle*{1.5}}
	\put(8,25){$b_{i, j}$}
	\put(13,28){\line(1,1.2){4}}
	\put(13,28){\line(1,-2.3){12}}
	
	\put(17,33){\circle*{1.5}}
	\put(19,32){$c_{i, j}$}
	\put(17,33){\line(1,0.9){7.3}}
	\put(17,33){\line(1,-4){8.3}}
	\end{picture}
	$$
	\caption{The graph $G$}
	\label{fig1}
\end{figure}



For clarity, we construct the graph  $G$ in a few steps. In the following,
vertices of different names are distinct.
\begin{enumerate}
	\item Add $p$ edges $x_iy_i$, $1 \le i \le p$.
	\item Add vertices $u$ and $z$, and let $u,z$ be  adjacent to $x_i,y_i$ for $1 \le i \le p$.
	\item For $1 \le i \le p$, 
	denote by $T_i$ the triangle induced by $\{x_i,y_i,z\}$. Inside $T_i$,
	add a path $(z_{i,1}, z_{i,2}, \ldots, z_{i, k},z)$, and let each $z_{i,j}$ be adjacent to both $x_i$ and $y_i$.
	\item For $1 \le j \le k$, $1 \le j \le k$,  denote by $T_{i,j}$ the triangle induced by $\{x_i, z_{i,j}, z_{i, j+1}\}$ (here we let $z_{i, k+1}=z$ for   $1 \le i \le p$).
	Inside $T_{i,j}$, add a triangle $(a_{i,j}, b_{i,j}, c_{i,j})$, and let 
	\subitem $a_{i,j}$ be adjacent to $x_i, z_{i,j}$, 
	\subitem  $b_{i,j}$ be adjacent to $x_i,z_{i,j+1}$, 
	\subitem $c_{i,j}$ be adjacent to $z_{i,j}$ and $z_{i,j+1}$.
\end{enumerate}
This completes the construction of the graph $G$.

So the   vertex set of $G$ is
$$V=\{u, z\} \cup \{x_i, y_i, z_{i,j}, a_{i,j}, b_{i,j},c_{i,j}: 1 \le i \le p, 1 \le j \le k\}.$$

Now we define a $4m$-list assignment $L$ of $G$. Let $A,B$ be two disjoint  sets of colours with $|A|=|B|=4m$. Let $L(u)=A$ and $L(z)=B$.
Let $${\cal S}=\left\{(X,Y): X \in {A \choose m}, Y \in {B \choose m}\right\}.$$
I.e., ${\cal S}$ is the family of pairs $(X,Y)$ such that $X$ is an $m$-subset of $A$ and $Y$ is an $m$-subset of $B$. Thus $|{\cal S}| = {4m \choose m}^2 =p$.
Order the elements of ${\cal S}$ arbitrarily, and denote by $(X_i,Y_i)$ the $i$th element of ${\cal S}$. Let $C,D$ be two disjoint sets, each of size $2m$ and $(C \cup D)  \cap (A \cup B)=\emptyset$. Let ${\cal T}$ be the family of $m$-subsets of $C$. So $|{\cal T}| = {2m \choose m}=k$. Order the elements of ${\cal T}$ arbitrarily, and denote by $Z_j$ the $j$th element of ${\cal T}$.
We define the lists of vertices inside $T_i$ as follows:
\begin{enumerate}
	\item   $L(x_i)=L(y_i)=L(z_{i,j})=C \cup X_i \cup Y_i$. 
	\item For $1 \le j \le k$,
	\subitem if $j$ is even, then 
	$L(a_{i,j})= Z_j \cup X_i \cup D$, 
	$L(b_{i,j})=Z_j \cup Y_i \cup D$,
	and  $L(c_{i,j})=X_i\cup Y_i \cup D$.
	\subitem If $j$ is odd, then 
	$L(a_{i,j})= Z_j \cup Y_i \cup D$, 
	$L(b_{i,j})=Z_j \cup X_i \cup D$,
	and $L(c_{i,j})=X_i\cup Y_i \cup D$.
\end{enumerate}
This completes the definition of $L$.  

Now we show that there is no $m$-fold $L$-colouring of $G$. Assume to the contrary that $\phi$ is an $m$-fold $L$-colouring of $G$. Then $(\phi(u), \phi(z)) \in {\cal S}$. Assume $(\phi(u), \phi(z))$ is the $i$th element of 
${\cal S}$. I.e., $\phi(u)=X_i$ and $\phi(z)=Y_i$. As $x_i,y_i$ are adjacent to $u$ and $z$, and $L(x_i)=L(y_i)= C \cup X_i \cup Y_i$, we conclude that $\phi(x_i) \cup \phi(y_i) = C$. As   $L(z_{i,j})=C \cup X_i \cup Y_i$ for $1 \le j \le k+1$ (recall that we let $z_{i,k+1}=z$), and $z_{i,j}$ is adjacent to both $x_i$ and $y_i$, we conclude that for $j=1,2,\ldots,k$, $\phi(z_{i,j}) \cup \phi(z_{i, j+1}) = X_i \cup Y_i$. As $z_{i,k+1} = z$ and $\phi(z)=Y_i$, we conclude that $\phi(z_{i,j}) = X_i$ if $j$ is even and $\phi(z_{i,j})=Y_i$ if $j$ is odd (observe that $k$ is even).




Now $\phi(x_i)$ is an $m$-subset of $C$, i.e., 
$\phi(x_i) \in {\cal T}$.
Assume $\phi(x_i)$ is the $j$th element of ${\cal T}$, i.e., $\phi(x_i)=Z_j$. 

Assume $j$ is even. As  $L(a_{i,j})= Z_j \cup X_i \cup D$, $a_{i,j}$ is adjacent to $x_i$ and $z_{i,j}$ and $\phi(z_{i,j})=X_i$, we conclude that $\phi(a_{i,j}) \subseteq D$. Similarly, we can prove that $\phi(b_{i,j}) \subseteq D$ and $\phi(c_{i,j}) \subseteq D$. As $|D|=2m$ and $\phi(a_{i,j}), \phi(b_{i,j}), \phi(c_{i,j})$ are pairwise disjoint $m$-sets, this is a contradiction. If $j$ is odd, we obtain a contradiction   similarly. This completes the proof of Theorem \ref{main}. \qed

The graph $G$ can be made a little smaller.  For example, we can add an edge connecting $u$ and $z$, and let $L(u)=L(z)$. In this way, we can choose $p={4m \choose 2m} {2m \choose m}$. Also we can add a triangle inside each triangle induced by $\{y_i,z_{i,j}, z_{i,j+1}\}$, and let $k = {2m \choose m}/2$.
However, the graph $G$ will still be a huge graph if  $m$ is large.

 